\chapter{Datos y balance}
\label{chap:datos}

\noindent Este capítulo describe los datos con los que se construye el análisis y la muestra sobre la que se estima el efecto de las cámaras. La primera sección presenta las fuentes, organizadas según el papel que cumple cada variable: el tratamiento, la siniestralidad, la infraestructura vial y el Metro. Para cada fuente se describe qué contiene, cómo se limpió y qué variables se construyen a partir de cada una. La segunda sección especifica cómo se construyen las unidades tratadas y los controles, y cuántas unidades quedan en cada paso. La tercera describe el modelo del \emph{propensity score} y las reglas del emparejamiento, y evalúa el soporte común y el balance de la muestra final. La última contiene las limitaciones de los datos.

\section{Fuentes y variables}

\subsection{Tratamiento}

\noindent La ubicación de las cámaras proviene del conjunto de datos \emph{Fotoc\'ivicas (Ubicación)} de la Secretaría de Seguridad Ciudadana (SSC, 2023), publicado en el portal de datos abiertos de la Ciudad de México. Para cada una de las 113 cámaras del programa reporta la vialidad, el cruce en el que se encuentra, el sentido de circulación que vigila y sus coordenadas. A diferencia del esquema de fotomultas, la ubicación de las cámaras de fotocívicas es pública, lo que permite identificar con precisión las zonas tratadas.\\

Las 113 cámaras se conservan sin limpieza. Tres pares de cámaras comparten exactamente la misma coordenada.\footnote{Son los registros 7 y 8, ambos descritos en el cruce de Insurgentes y Eje 1 Norte, en Buenavista; 5 y 25, descritos en Universidad y Miguel Ángel de Quevedo y en Eje 5 Sur y Anillo Periférico Oriente, y 52 y 53, descritos sobre la Autopista Urbana Norte a la altura de Constituyentes y de Ejército Nacional. En los dos últimos pares, la coordenada coincide con el cruce de la primera cámara: está a 40 metros de la estación Miguel Ángel de Quevedo del Metro y a 200 metros del cruce de Constituyentes con el Periférico. Los cruces descritos para las cámaras 25 y 53 están a 13 y 4 kilómetros de la coordenada, respectivamente.} En el primer par se trata de dos cámaras en el mismo cruce, una para cada sentido. En los otros dos, el cruce que describe la fuente para una de las cámaras está a varios kilómetros de la coordenada registrada, es decir, la fuente le asignó la coordenada de la otra cámara. Como no hay forma de conocer su ubicación exacta, cada par se trata como cualquier otro par de cámaras cercanas: se agrupan en una sola unidad tratada (sección~\ref{sec:unidades}).\\

A partir de estas ubicaciones se define el indicador de tratamiento $T_i$, que vale uno si el área $i$ contiene al menos una cámara y cero si es un control.

\subsection{Siniestralidad}

\noindent Los incidentes viales provienen del conjunto de datos \emph{Incidentes viales reportados por C5} del Centro de Comando, Control, Cómputo, Comunicaciones y Contacto Ciudadano de la Ciudad de México (C5, 2024). Cada registro corresponde a un folio de atención, con su fecha, hora, coordenadas, tipo de incidente y la clasificación con la que se cerró. Se prefiere esta fuente frente a los registros de la SSC por dos razones. La primera es su cobertura temporal: los datos de la SSC solo tienen coordenadas geográficas a partir de junio de 2018, lo que impide construir tres años previos al programa, mientras que el C5 reporta incidentes desde diciembre de 2013. La segunda es su alcance: la base de la SSC se limita a incidentes con lesionados o fallecidos, mientras que el C5 integra llamadas al 911, botones de auxilio y videovigilancia, e incluye incidentes menores. Excluirlos impediría detectar un efecto del programa sobre la frecuencia de los incidentes, y no solo sobre su gravedad.\\

Los datos contienen 2,115,080 folios entre el 31 de diciembre de 2013 y el 29 de febrero de 2024. Para la limpieza se hizo lo siguiente:

\begin{enumerate}
    \item Se conservan los folios con código de cierre A (afirmativo), es decir, aquellos en los que una unidad de emergencia llegó al lugar y confirmó el incidente. Quedan 866,230.
    \item Se recorta la ventana del análisis, del 22 de abril de 2016 al 21 de abril de 2022. Quedan 433,235.
    \item Se conservan los incidentes tipificados como \emph{Accidente}, \emph{Cadáver} o \emph{Lesionado}, sin la clasificación de falsa alarma, para descartar llamadas que no corresponden a hechos de tránsito. Quedan 432,707.
    \item Se eliminan cuatro folios repetidos y cinco registros con coordenadas en cero. Quedan 432,698.
\end{enumerate}

Los incidentes se clasifican en tres niveles mutuamente excluyentes según su gravedad: fatales (FCS), si se registró un cadáver en el lugar, con 1,957 incidentes; con lesionados (PIC), si la clasificación de la alarma fue de urgencias médicas, con 191,676, y menores (MIN), sin lesionados ni fallecidos, con 239,065.\\

\textbf{Variable de resultado}. La variable de resultado $Y_{it}$ es el número de incidentes en el área $i$ durante el periodo $t$, tanto el total como para cada nivel de severidad. Los periodos van del día 22 de un mes al 21 del siguiente, de modo que el corte del 22 de abril de 2019 coincide con el inicio de un periodo, y la ventana tiene 36 periodos antes del programa y 36 después. Un incidente que cae en el área compartida por dos o más áreas tratadas se reparte en partes iguales entre ellas, como se detalla en la sección~\ref{sec:unidades}.\\

\textbf{Covariables de siniestralidad}. La misma fuente da lugar a cuatro covariables, que describen la siniestralidad de cada área en los tres años previos al programa. Se consideran los años hacia atrás a partir del 22 de abril de 2019: el año $-1$ va del 22 de abril de 2018 al 21 de abril de 2019, y así sucesivamente.

\begin{itemize}
    \item \textbf{Nivel de incidentes} (\texttt{nivel\_incidentes}): promedio mensual de incidentes de todos los niveles en los tres años previos.
    \item \textbf{Proporción con lesionados} (\texttt{pct\_lesionados}): fracción de esos incidentes que tuvo consecuencias personales; incluye tanto PIC como FCS.
    \item \textbf{Tendencia} (\texttt{tendencia}): promedio mensual de incidentes del año $-1$ menos el del año $-3$. Un valor positivo indica que los incidentes del área venían en aumento antes del programa, y uno negativo, que venían a la baja.
    \item \textbf{Algún incidente fatal} (\texttt{hubo\_fcs}): indicador de que el área registró al menos un incidente fatal en los tres años previos. Se usa un indicador y no un conteo porque los incidentes fatales son poco frecuentes: en un año dado, entre el 58\% y el 84\% de las áreas no registra ninguno.
\end{itemize}

Con estas covariables se representan los factores que consideró la Secretaría de Movilidad al reubicar las cámaras, cuando se priorizaron los lugares donde los incidentes eran más frecuentes y más graves (SEMOVI, 2019). La tendencia, además, permite que el emparejamiento compare áreas cuya siniestralidad venía evolucionando de forma parecida antes del programa, lo que sirve directamente al supuesto de tendencias paralelas del modelo de diferencias en diferencias.

\subsection{Infraestructura vial}

\noindent Las características de la red vial provienen del conjunto de datos \emph{Vialidades primarias de la Ciudad de México} de la Secretaría de Movilidad (SEMOVI, 2023). Contiene los tramos de las vías primarias y de acceso controlado de la ciudad, con el número de carriles, el sentido de circulación, el tipo de vía y el nivel en el que corre cada tramo: a nivel de calle, elevado o subterráneo. En total se identifican 149 vialidades: 135 con tramos de vía primaria y 21 con tramos de acceso controlado, ya que siete tienen tramos de ambos tipos. La misma red sirve para sortear los centros de los controles (sección~\ref{sec:controles}).\\

A partir de ella se construyen siete covariables, ponderadas por los metros de vía que caen dentro de cada área:

\begin{itemize}
    \item \textbf{Longitud vial} (\texttt{road\_length\_m}): metros de vía dentro del área.
    \item \textbf{Número de vialidades} (\texttt{n\_vialidades}): vialidades distintas que cruzan el área.
    \item \textbf{Carriles ponderados} (\texttt{carriles\_ponderados}): número promedio de carriles, ponderado por los metros de cada tramo.
    \item \textbf{Acceso controlado} (\texttt{pct\_acceso\_controlado}): fracción de los metros de vía que corresponde a vías de acceso controlado.
    \item \textbf{Doble sentido} (\texttt{pct\_doble\_sentido}): fracción de los metros de vía en tramos de doble sentido.
    \item \textbf{Desnivel} (\texttt{pct\_desnivel}): fracción de los metros de vía en tramos elevados o bajo el suelo.
    \item \textbf{Número de niveles} (\texttt{n\_niveles}): niveles distintos que se encuentran dentro del área.
\end{itemize}

Las covariables se ponderan por longitud porque así describen la composición de toda el área, y no solo la vía más grande que la cruza: un círculo atravesado por unos metros de una avenida de seis carriles no se parece a uno recorrido por ella de lado a lado. No se incluye la fracción de vía primaria porque, al contener la capa solo vías primarias y de acceso controlado, es igual a uno menos la fracción de acceso controlado. Las dos últimas covariables miden la ambigüedad vertical de cada área, que se discute en la sección~\ref{sec:limitaciones}.

\subsection{Metro}

\noindent La información del Metro combina dos fuentes de la Secretaría de Movilidad. La primera es la capa geográfica \emph{Ubicación de líneas y estaciones del Sistema de Transporte Colectivo Metro} (SEMOVI, 2026), con la ubicación de las estaciones. Tiene 195 registros para 163 estaciones, porque las estaciones de transbordo tienen un registro por línea. La segunda es la \emph{Afluencia diaria del Metro CDMX} (SEMOVI, 2025), que reporta el número de usuarios por día, línea y estación desde 2010.\\

Para poder unir ambas fuentes fue necesario corregir el nombre de dos estaciones: la afluencia usa el nombre anterior de Niños Héroes, y la capa geográfica escribe ``Mixhiuca'' en lugar de Mixiuhca. Con estas correcciones, todos los pares de línea y estación de la afluencia encuentran su ubicación.\\

De estas fuentes sale una sola covariable:

\begin{itemize}
    \item \textbf{Afluencia del Metro} (\texttt{afluencia\_nivel}): promedio mensual de usuarios de las estaciones ubicadas dentro del área, en los tres años previos al programa. Vale cero si el área no contiene ninguna estación.
\end{itemize}

La afluencia funciona como un indicador de la actividad y del tránsito de personas en la zona. Basta con un solo promedio porque la afluencia de una estación casi no cambia de un año a otro: lo que distingue a las áreas es qué estación tienen. A 300 metros, el 32\% de las áreas tratadas tiene alguna estación del Metro dentro, mientras que solo el 13\% de los controles la tiene. El cero representa una zona sin Metro, no falta de información. Se prefirió esta medida a la distancia a la estación más cercana, que obligaría a suponer que una estación influye en la zona incluso cuando queda fuera del área.

\subsection{Resumen de variables}

\noindent El cuadro~\ref{tab:variables} resume las variables del análisis. Las doce covariables se miden entre el 22 de abril de 2016 y el 21 de abril de 2019, y cada una tiene un solo valor por área, que permanece fijo tras la implementación del programa. Las doce entran al modelo del \emph{propensity score}. En la especificación principal~\eqref{eq:did}, las covariables se interactúan con el periodo posterior, excepto las de siniestralidad, que se excluyen por las razones expuestas en la sección~\ref{sec:especificacion}.

\begin{table}[htbp]
    \centering
    \caption{Variables del análisis}
    \label{tab:variables}
    \footnotesize
    \begin{tabularx}{\textwidth}{l l >{\raggedright\arraybackslash}X}
        \toprule
        Variable & Fuente & Definición \\
        \midrule
        \multicolumn{3}{l}{\textit{Tratamiento y resultado}} \\
        $T_i$ & SSC & Indicador de que el área contiene al menos una cámara \\
        $Y_{it}$ & C5 & Incidentes en el área $i$ durante el periodo $t$, en total y por nivel \\
        \midrule
        \multicolumn{3}{l}{\textit{Siniestralidad previa}} \\
        \texttt{nivel\_incidentes} & C5 & Incidentes al mes, promedio de los tres años previos \\
        \texttt{pct\_lesionados} & C5 & Fracción de incidentes con lesionados o fallecidos \\
        \texttt{tendencia} & C5 & Incidentes al mes del año $-1$ menos los del año $-3$ \\
        \texttt{hubo\_fcs} & C5 & Al menos un incidente fatal en los tres años previos \\
        \midrule
        \multicolumn{3}{l}{\textit{Infraestructura vial}} \\
        \texttt{road\_length\_m} & SEMOVI & Metros de vía dentro del área \\
        \texttt{n\_vialidades} & SEMOVI & Vialidades distintas que cruzan el área \\
        \texttt{carriles\_ponderados} & SEMOVI & Carriles promedio, ponderados por metros de vía \\
        \texttt{pct\_acceso\_controlado} & SEMOVI & Fracción de los metros en vías de acceso controlado \\
        \texttt{pct\_doble\_sentido} & SEMOVI & Fracción de los metros en vías de doble sentido \\
        \texttt{pct\_desnivel} & SEMOVI & Fracción de los metros en tramos elevados o bajo el suelo \\
        \texttt{n\_niveles} & SEMOVI & Niveles distintos dentro del área \\
        \midrule
        \multicolumn{3}{l}{\textit{Movilidad}} \\
        \texttt{afluencia\_nivel} & SEMOVI & Usuarios al mes de las estaciones del Metro dentro del área; cero si no hay estación \\
        \bottomrule
    \end{tabularx}
\end{table}

\section{Unidades de análisis y muestra}

\noindent La unidad de análisis es un área circular de radio $r$, con $r \in \{100, 150, 200, 250, 300\}$ metros, observada en cada uno de los 72 periodos de la ventana. Tanto las áreas tratadas como las de control son círculos del mismo radio, lo que hace posible compararlas entre sí. A continuación se especifica cómo se definen los centros de los círculos de ambos tipos de unidades.

\subsection{Unidades tratadas}
\label{sec:unidades}

\noindent Algunas cámaras están tan cerca entre sí que sus círculos prácticamente coinciden, y tratarlas por separado contaría dos veces la misma zona. Por ello, dos cámaras se agrupan en una sola unidad tratada cuando sus círculos de 300 metros comparten más del 60\% de su área, lo que equivale a que sus ubicaciones estén a menos de 192 metros. El centro de cada unidad es el promedio de las ubicaciones de sus cámaras. Tras la construcción, las 113 cámaras quedan en 99 unidades tratadas: 87 con una sola cámara, 10 con dos y 2 con tres.\\

El conjunto de unidades tratadas se construye una sola vez, con el radio de 300 metros, y se mantiene igual en el resto de los radios. Así, siempre se habla de los mismos centros, y lo único que cambia entre estimaciones es el radio.\\

Dos condiciones acotan el umbral del 60\%. Por un lado, con un umbral de 50\% se agruparían dos cámaras separadas por 204 metros, cuyo centro queda a 102 metros de cada una: a 100 metros, esa unidad tratada no contendría ninguna cámara. Con 60\%, la cámara más alejada del centro de su unidad está a 81 metros. Por otro lado, con un umbral de 70\% se agrupan menos cámaras y aparecen círculos traslapados desde los 100 metros, mientras que con 60\% no hay ningún traslape en ese radio.\footnote{Ningún umbral evita todos los casos límite. Las unidades 32 y 77 son dos cámaras sobre el mismo tramo del segundo piso del Periférico, con sus centros a 205 metros, y quedaron separadas por 13 metros. Con un umbral de 50\%, que sí las agrupa, el efecto estimado a 300 metros en la especificación principal cambia a lo más 0.18 errores estándar en los cuatro niveles de gravedad.}\\

\textbf{Traslape entre unidades tratadas}. Con esta regla, los círculos de unidades distintas pueden traslaparse parcialmente en los radios mayores. A 100 metros no hay ningún traslape; a 300 metros, 36 unidades comparten parte de su área con alguna otra, y en esas áreas compartidas cae el 5.5\% de los incidentes. Cada incidente en un área compartida se reparte en partes iguales entre las unidades que lo contienen: cuenta $1/2$ en cada una si cae en dos círculos, $1/3$ si cae en tres, y así sucesivamente. Con este reparto, la suma de los incidentes de todas las unidades coincide con el número real de incidentes, 44,153 a 300 metros, en lugar de los 46,730 que resultarían de contar dos veces los compartidos.\\

El reparto resuelve el conteo de los incidentes, pero implica que los incidentes de unidades traslapadas dejan de ser independientes. Por ello, las unidades tratadas conectadas por algún traslape forman una \emph{componente}, que es la que se usa para agrupar los errores estándar y para remuestrear en el \emph{bootstrap}. Es importante distinguir que una componente no es una unidad de análisis: cada unidad conserva su propio círculo y su propio emparejamiento, y las componentes cambian con el radio. La figura~\ref{fig:componentes} muestra un ejemplo: las mismas cinco unidades tratadas forman cinco componentes a 100 metros, cuatro a 200 y una sola a 300. Entre las 93 unidades de la muestra final hay 93 componentes a 100 metros, 86 a 200 y 72 a 300.

\begin{figure}[htbp]
    \centering
    \includegraphics[width=\linewidth]{Imagenes/componente-tratado-100-200-300.pdf}
    \caption{Las mismas cinco unidades tratadas en tres radios}
    \label{fig:componentes}
    \begin{minipage}{0.95\linewidth}
        \footnotesize \textit{Nota:} Cada círculo es una unidad tratada, centrada en el promedio de sus cámaras. Las líneas rojas unen las unidades cuyos círculos se traslapan; las unidades conectadas forman una componente. En gris, la red de vías primarias y de acceso controlado. La escala es la misma en los tres paneles.
    \end{minipage}
\end{figure}

\subsection{Controles}
\label{sec:controles}

\noindent La construcción de los controles se hace sobre la misma red vial, considerando tanto las vías primarias como las de acceso controlado. Se colocan 20,000 puntos al azar, con probabilidad proporcional a la longitud de cada tramo, de modo que cada metro de vía tiene la misma probabilidad de recibir un punto. Se eligió esta construcción, frente a usar los cruces de la red o una malla, porque es la que menos decisiones arbitrarias requiere: no hay que definir qué cuenta como un cruce ni dónde empieza la malla, y cubre por igual los tramos y los cruces.\\

Una vez sorteados los puntos, se eliminan los que quedan a menos de 600 metros del centro de alguna unidad tratada, es decir, a menos de dos veces el radio más grande. Con esto, ningún control comparte área con una unidad tratada en ninguno de los cinco radios. Quedan 14,656 posibles controles, alrededor de 148 por cada unidad tratada, y este conjunto también se mantiene fijo para todos los radios.\\

Inicialmente se permite que los candidatos se traslapen entre sí: cada candidato representa el contrafactual de su propia zona, y exigir que no se traslaparan reduciría el conjunto a unos cientos de puntos y empeoraría el emparejamiento. Por la misma razón, los incidentes de los candidatos cuentan completos y no se reparten. La restricción se impone sobre los controles seleccionados: en el emparejamiento, ningún control puede quedar a menos de dos radios de otro control ya elegido, de modo que los controles de la muestra final no se traslapan entre sí y cada uno forma su propia componente.

\subsection{Muestra final}

\noindent Dado que las covariables se calculan dentro del círculo, el emparejamiento se hace por separado en cada radio. En cada radio, algunas unidades tratadas no encuentran un control suficientemente parecido (sección~\ref{sec:emparejamiento}): tres a 100 metros, una a 150, ninguna a 200, dos a 250 y una a 300. Para que las estimaciones de los cinco radios describan exactamente las mismas zonas, la muestra final se restringe a las 93 unidades tratadas que encuentran pareja en los cinco radios. Los controles sí se eligen de nuevo en cada radio. En resumen:

\begin{itemize}
    \item 113 cámaras forman 99 unidades tratadas;
    \item 20,000 puntos sorteados dejan 14,656 candidatos a control;
    \item 93 unidades tratadas encuentran pareja en los cinco radios, y cada una se empareja con un control.
\end{itemize}

Cada radio se estima, entonces, sobre 186 áreas observadas en 72 periodos. El efecto estimado es el efecto promedio sobre estas 93 unidades tratadas, y no sobre las 99; la sección~\ref{sec:excluidas} describe a las seis que quedan fuera.

\section{Emparejamiento y balance}
\label{sec:emparejamiento}

\subsection{\emph{Propensity score}}

\noindent El \emph{propensity score} se estima con una regresión logística del indicador de tratamiento sobre las doce covariables, con las 99 unidades tratadas y los 14,656 candidatos. La estimación no usa penalización, porque el objetivo no es predecir bien en una muestra que no se observa, sino balancear las covariables en la muestra con la que se trabaja; una regresión penalizada encoge los coeficientes y deja a propósito parte del desbalance que se quiere corregir. Las covariables se estandarizan antes de estimar. La regresión logística no depende de la escala de las variables, pero estandarizarlas hace que los coeficientes sean comparables entre sí y que el algoritmo de estimación converja. El modelo se estima con el método de Newton, que es el procedimiento estándar para la regresión logística, y converge en siete iteraciones. Con el algoritmo por defecto de la biblioteca utilizada, L-BFGS de \emph{scikit-learn}, el modelo no converge si las covariables no se estandarizan y, aun estandarizadas, se detiene antes de llegar al óptimo, con diferencias que alcanzan para cambiar hasta 18 de los 99 pares.

\subsection{Reglas del emparejamiento}

\noindent Cada unidad tratada se empareja con un control, sin reemplazo, de acuerdo con cuatro reglas:

\begin{itemize}
    \item \textbf{Caliper}. Un control solo es elegible si su logit del \emph{propensity score} está a menos de 0.2 desviaciones estándar del de la unidad tratada, que es el ancho recomendado en la literatura (Austin, 2011). Una unidad tratada sin controles dentro del caliper se queda sin pareja.
    \item \textbf{Desempate por distancia de Mahalanobis}. Dentro del caliper suele haber cientos de controles con un \emph{propensity score} casi idéntico; la diferencia mediana entre ellos es de 0.0002. Si se eligiera entre ellos por \emph{propensity score} sería casi como elegir al azar, por lo que se elige el más cercano en todas las covariables a la vez, con la distancia de Mahalanobis (Rosenbaum y Rubin, 1985). Esto reduce la diferencia estandarizada promedio de 0.093 a 0.070 sin perder ningún par.
    \item \textbf{Orden}. Las unidades tratadas eligen control de mayor a menor \emph{propensity score}, de modo que las más difíciles de emparejar eligen primero, cuando hay más controles disponibles.
    \item \textbf{Separación entre controles}. Un control no es elegible si queda a menos de dos radios de otro control ya elegido.
\end{itemize}

El emparejamiento se rehace en cada radio con las covariables medidas a ese radio. Reutilizar en todos los radios los pares formados a 300 metros deteriora el balance, porque un control que se parece a su unidad tratada en un círculo de 300 metros no necesariamente se le parece en uno de 100.

\subsection{Soporte común}

\noindent El emparejamiento solo es defendible donde las distribuciones del \emph{propensity score} de tratadas y controles se traslapan; de hacerlo fuera de esa zona, la comparación sería difícil de sostener. A 300 metros, el \emph{propensity score} de las tratadas va de 0.001 a 0.140, con mediana de 0.011, y el de los candidatos va de casi cero a 0.139, con mediana de 0.004. Solo una unidad tratada queda por encima del candidato más alto. Las tratadas se concentran en valores más altos que los candidatos, como se esperaría si las cámaras se ubicaron en zonas más riesgosas, pero incluso en el quintil más alto de las tratadas hay 740 candidatos en el mismo rango. La figura~\ref{fig:soporte} muestra las dos distribuciones antes y después del emparejamiento.

\begin{figure}[htbp]
    \centering
    \includegraphics[width=\linewidth]{Imagenes/soporte-comun-300.pdf}
    \caption{Distribución del \emph{propensity score} a 300 metros}
    \label{fig:soporte}
    \begin{minipage}{0.95\linewidth}
        \footnotesize \textit{Nota:} El panel (a) compara las 99 unidades tratadas con los 14,656 candidatos a control; el panel (b), las 93 unidades tratadas de la muestra final con sus controles emparejados. El eje horizontal está en escala logarítmica porque la mayoría de los valores se concentra cerca de cero.
    \end{minipage}
\end{figure}

\subsection{Balance}

\noindent El balance se mide con la diferencia de medias estandarizada (SMD) de cada covariable:
\[
    \text{SMD} = \frac{\bar{X}_T - \bar{X}_C}{\sqrt{(s_T^2 + s_C^2)/2}},
\]
donde $\bar{X}$ y $s^2$ son la media y la varianza de la covariable en las tratadas ($T$) y en los controles ($C$). A diferencia de una prueba de hipótesis, no depende del tamaño de la muestra. Se considera que una covariable está balanceada cuando el valor absoluto de su SMD es menor a 0.1 (Austin, 2009).\\

El cuadro~\ref{tab:balance} reporta las medias y las SMD a 300 metros. Antes del emparejamiento, las tratadas tienen más incidentes, más incidentes fatales, más carriles y más afluencia del Metro que los candidatos, lo que confirma que las cámaras no se ubicaron al azar: la SMD promedio en valor absoluto es de 0.266, y diez de las doce covariables superan el umbral. Después del emparejamiento, la SMD promedio baja a 0.065 y solo tres covariables lo superan: la longitud vial (0.120), la proporción con lesionados (0.115) y la afluencia del Metro (0.141). En los demás radios, la SMD promedio después del emparejamiento va de 0.045 a 0.055, y entre una y dos covariables superan el umbral; el valor más alto es de 0.268, a 200 metros.

\begin{table}[htbp]
    \centering
    \caption{Balance de covariables a 300 metros}
    \label{tab:balance}
    \footnotesize
    \setlength{\tabcolsep}{4pt}
    \begin{tabular}{l r r r r r}
        \toprule
        & \multicolumn{3}{c}{Media} & \multicolumn{2}{c}{SMD} \\
        \cmidrule(lr){2-4} \cmidrule(lr){5-6}
        Covariable & Tratadas & Candidatos & Controles & Antes & Después \\
        \midrule
        \texttt{nivel\_incidentes} & 6.22 & 4.57 & 6.02 & 0.634 & 0.070 \\
        \texttt{pct\_lesionados} & 0.40 & 0.40 & 0.39 & 0.000 & 0.115 \\
        \texttt{tendencia} & 0.13 & $-0.08$ & 0.20 & 0.168 & $-0.051$ \\
        \texttt{hubo\_fcs} & 0.70 & 0.41 & 0.72 & 0.647 & $-0.047$ \\
        \texttt{road\_length\_m} & 2,686 & 2,405 & 2,533 & 0.238 & 0.120 \\
        \texttt{n\_vialidades} & 2.15 & 1.97 & 2.09 & 0.196 & 0.066 \\
        \texttt{carriles\_ponderados} & 3.05 & 2.75 & 3.04 & 0.495 & 0.015 \\
        \texttt{pct\_acceso\_controlado} & 0.36 & 0.30 & 0.35 & 0.161 & 0.013 \\
        \texttt{pct\_doble\_sentido} & 0.11 & 0.07 & 0.10 & 0.179 & 0.047 \\
        \texttt{pct\_desnivel} & 0.10 & 0.10 & 0.09 & $-0.020$ & 0.051 \\
        \texttt{n\_niveles} & 1.59 & 1.51 & 1.56 & 0.128 & 0.042 \\
        \texttt{afluencia\_nivel} & 237,667 & 107,239 & 173,912 & 0.322 & 0.141 \\
        \bottomrule
    \end{tabular}
    \begin{minipage}{0.95\linewidth}
        \vspace{0.4em}
        \footnotesize \textit{Nota:} Las columnas \emph{Tratadas} y \emph{Controles} corresponden a las 93 unidades tratadas de la muestra final y a sus controles emparejados; \emph{Candidatos}, a los 14,656 candidatos a control. La SMD antes del emparejamiento compara las 99 unidades tratadas con los candidatos.
    \end{minipage}
\end{table}

Las diferencias que quedan se atienden de dos formas en la estimación. Los efectos fijos por unidad absorben cualquier diferencia de nivel entre tratadas y controles que no cambie en el tiempo, de modo que una diferencia en la longitud vial o en la afluencia del Metro no sesga el efecto por sí misma. Lo que sí podría sesgarlo es que esas diferencias hicieran evolucionar distinto a ambos grupos después de 2019, y para eso la especificación principal interactúa las covariables de infraestructura vial y la afluencia del Metro con el periodo posterior. La proporción con lesionados no se incluye en esa interacción, por las razones expuestas en la sección~\ref{sec:especificacion}; que su diferencia no se traduzca en trayectorias distintas es lo que evalúa el \emph{event study} del capítulo 4.

\subsection{Unidades excluidas}
\label{sec:excluidas}

\noindent Las seis unidades tratadas que quedan fuera de la muestra final no son un descarte al azar: son las zonas más peligrosas de la red. Antes del programa registraban 13.5 incidentes al mes, 2.2 veces el promedio de las 93 analizadas. Las seis habían tenido al menos un incidente fatal, contra el 70\% de las demás, tienen más carriles y más vía de acceso controlado, triplican la afluencia del Metro y su siniestralidad venía creciendo más rápido (figura~\ref{fig:excluidas}). No se parecen a ningún candidato a control, y por eso el emparejamiento no les encuentra pareja en al menos uno de los radios. La consecuencia es que los resultados de esta tesis no se extienden a ellas: si las cámaras tuvieran un efecto mayor justo donde el riesgo es más alto, este análisis no podría detectarlo.

\begin{figure}[htbp]
    \centering
    \includegraphics[width=\linewidth]{Imagenes/unidades-excluidas.pdf}
    \caption{Las seis unidades excluidas frente a las 93 analizadas}
    \label{fig:excluidas}
    \begin{minipage}{0.95\linewidth}
        \footnotesize \textit{Nota:} Cada punto es el promedio del grupo a 300 metros dividido entre el de las 93 unidades tratadas analizadas, cuyo valor aparece entre paréntesis. La línea punteada marca el valor de las 93 tratadas.
    \end{minipage}
\end{figure}

\section{Limitaciones de los datos}
\label{sec:limitaciones}

\begin{itemize}
    \item \textbf{No se observa el tránsito vehicular}. No se cuenta con información de cuántos vehículos circulaban dentro de la ciudad ni de las velocidades a las que lo hacían. Las áreas circulares mitigan en parte el desvío del tránsito hacia vías alternas, pero no es posible confirmarlo, y tampoco se puede saber si las cámaras modificaron las velocidades. Los únicos datos públicos de aforo vehicular a los que se tuvo acceso son los de las casetas de acceso a la ciudad, pero forman una sola serie para toda la ciudad, que los efectos fijos por periodo ya absorben.
    \item \textbf{No se sabe en qué nivel de la vialidad ocurrió cada incidente}. Si bien el C5 reporta las coordenadas de cada incidente, no especifica en qué nivel ocurrió, de modo que un incidente en un segundo piso no se distingue de uno a nivel de calle. El 42\% de las cámaras tiene vías de más de un nivel dentro de su círculo de 300 metros, y en esas áreas se cuentan incidentes sobre los que la cámara no podía influir, lo que diluye el efecto estimado hacia cero. Las covariables de desnivel y número de niveles buscan controlar esta ambigüedad, pero no la eliminan.
    \item \textbf{Dos cámaras no tienen su ubicación real}. Las cámaras 25 y 53 tienen la coordenada de otra cámara, y los cruces que describe la fuente para ellas están a varios kilómetros. Según la red vial, en ambos cruces hay otras cámaras del programa a menos de 100 metros, por lo que esas zonas probablemente ya forman parte de alguna unidad tratada, pero no es posible confirmarlo.
    \item \textbf{Solo se observan los incidentes reportados y confirmados}. Los incidentes menores que no se reportaron al C5, o en los que ninguna unidad de emergencia confirmó el hecho, no aparecen en los datos. Si el registro hubiera cambiado de forma distinta en las zonas tratadas y en sus controles después del programa, los resultados estarían sesgados.
\end{itemize}
